\section{The Query That Stopped Working}
\label{sec:ch10-the-query-that-stopped-working}
\index{federated query!the first real one|(}%

This is the query chapters~\ref{ch:hotchocolate-quickly}
to~\ref{ch:schema-design-that-survives-change} opened with, in its current
spelling:

\begin{minted}{graphql}
{
  browseProducts(first: 3) {
    nodes {
      title
      price { amount currency }
      availableQuantity
      averageRating
      reviews(first: 2) {
        totalCount
        nodes { rating author { displayName } }
      }
    }
  }
}
\end{minted}

It has not answered since
chapter~\ref{ch:the-first-cut-extracting-catalog}, and the two failures are
worth seeing side by side. Both subgraphs reject it at validation, with a 400
and no execution at all:

\begin{minted}[fontsize=\footnotesize]{text}
catalog (5101), 4 errors
  The field `price` does not exist on the type `Product`.
  The field `availableQuantity` does not exist on the type `Product`.
  The field `averageRating` does not exist on the type `Product`.
  The field `reviews` does not exist on the type `Product`.

mosaic (5100), 1 error
  The field `browseProducts` does not exist on the type `Query`.
\end{minted}

Catalog has the products and none of the four fields hanging off them. Mosaic
has all four and no way to list a product.

Send it to the router:

\begin{minted}[fontsize=\footnotesize]{json}
{
  "title": "Beech Cutting Board",
  "price": { "amount": 69, "currency": "EUR" },
  "availableQuantity": 61,
  "averageRating": 4.166666666666667,
  "reviews": {
    "totalCount": 6,
    "nodes": [
      { "rating": 5, "author": { "displayName": "Katharina Brandt" } },
      { "rating": 4, "author": { "displayName": "Tomas Novak" } }
    ]
  }
}
\end{minted}

That is the first of the three products, and it is the first answer about
Mosaic that crossed a service boundary without me carrying it.
Chapter~\ref{ch:how-a-federated-query-actually-runs} watched a router assemble
one, but that was two subgraphs built to be watched. This is the storefront.
\code{title} came out of Catalog's database. Everything else came out of
Mosaic's, hung off a \code{Product} class that holds an identifier and nothing
else.

\index{federation!the payoff, stated plainly}%
I want to be precise about what is impressive here and what is not. The data is
dull: it is the same twenty-five products and hundred and twenty reviews the
book has been carrying since chapter~\ref{ch:hotchocolate-quickly}. What
changed is that no process in the system has all of it, and the client did not
have to know that. There is no aggregation code anywhere in Mosaic. Nobody
wrote a class that calls Catalog and then calls Pricing and stitches the
results. The stitching is a consequence of two schemas agreeing on what a
\code{Product} is, and the rest is a routing table.

\subsection*{What each service did for it}

\index{request timeline!through a router}%
Chapter~\ref{ch:the-life-of-a-request} built a per-request timeline and
chapter~\ref{ch:data-without-the-n-plus-1} added a SQL counter to it. Both are
still running. Mosaic's console, for one warm request:

\begin{minted}[fontsize=\footnotesize,breaklines]{text}
parse - validate - compile - coerce 0.056ms execute 7.802ms total 7.895ms (document cache hit, operation cache hit, 19 resolvers, 6 SQL)
Service lookups this request: 5
\end{minted}

Nineteen resolvers and six statements for its half. Both caches hit, which they
would: the router sends the same document every time, which is
section~\ref{sec:ch10-the-plan-on-a-real-graph}'s subject from the other end.
The first request of a session reports both caches missed and 93.9 ms of
execute, so chapter~\ref{ch:the-life-of-a-request}'s warning about cold first
measurements survives the move to a federated graph intact.

\index{observability!lost in the extraction}%
Catalog's console reports nothing at all. Not a smaller number: nothing. It has
no request timeline, no resolver count and no lookup counter, because
chapter~\ref{ch:the-life-of-a-request} built all three inside \code{Mosaic.Api}
and chapter~\ref{ch:the-first-cut-extracting-catalog}'s extraction left them
there. I grepped the whole of \code{src/Mosaic.Catalog} for the four
diagnostic types and got one hit, which is a comment in
\code{CatalogService.cs} saying the counter is gone.

That is worth more than a footnote. Half of every federated query in this
system is now unmeasured, and the extraction that caused it did not warn me,
because instrumentation is not part of a schema and no gate checks for it. The
service still works. It is just quiet. This is the ordinary way observability
degrades when a service splits, and
chapter~\ref{ch:observability} is where it gets fixed properly with tracing that
spans both.

\subsection*{Verify it in Postman}

\index{Postman!against the router}%
The collection for this chapter is \code{mosaic-router.postman\_collection.json},
and unlike chapter~\ref{ch:the-first-cut-extracting-catalog}'s it talks to one
address. Seven requests, twenty-three assertions, all against port 3002.

The first one sends the query above and asserts what makes it interesting rather
than what makes it pass:

\begin{minted}[fontsize=\footnotesize]{javascript}
pm.test("every product carries fields from both subgraphs", function () {
    for (const node of nodes) {
        pm.expect(node.title, "title").to.be.a("string").and.not.empty;
        pm.expect(node.price.amount, "price.amount").to.be.a("number");
        pm.expect(node.price.currency, "price.currency").to.eql("EUR");
        pm.expect(node.availableQuantity, "availableQuantity").to.be.a("number");
        pm.expect(node.reviews.totalCount, "reviews.totalCount").to.be.a("number");
    }
});
\end{minted}

\index{verification!asserting the seam}%
An assertion on the product count would pass against Catalog alone. This one
cannot: every field after \code{title} is on the far side of the seam, so the
test fails if the second hop stops happening, and it fails for every product
rather than for whichever one happened to be first. The verification script runs
it on every pass, and its line in the summary says what it is for:

\begin{minted}{text}
[ok]   the router answers the query neither subgraph can
\end{minted}

Two subgraphs, one answer, and so far nothing about how. The router will tell
you how if you ask it, and it turns out to have rewritten the question on the
way in.

\index{federated query!the first real one|)}%
